\section{Conclusion}

In this paper, we introduced \textbf{TURRET OS}, a multi-layered counter-espionage framework that unifies cross-format document metadata harvest, a W3C-PROV provenance knowledge graph, a boolean rule AST engine, a spatio-temporal GNN, and Ed25519-signed Merkle evidence packaging. Evaluated across 141,725 activity records (D3 TSEC corpus) and a 90-day lab deployment pilot (D5 synthetic tenant), TURRET OS achieves an uncalibrated ROC-AUC of 0.9983, a calibrated best F1-score of 0.9747 (at threshold 0.37), an F1-score of 0.8813 at a 0.5\% FPR operating point, an MCC of 0.9732, and an operational time-to-detect of $<0.1$ minutes. Under adversarial red-team evaluations, TURRET OS rescues single-layer detection degradation under OOD credential-phishing attacks from 0.7363 to 0.9572 ROC-AUC. Furthermore, TURRET OS delivers 100.0\% coverage across all eight ISO/IEC 27043 digital forensic readiness attributes, providing security operations centers with an automated bridge from endpoint telemetry to court-admissible forensic evidence.

\textbf{Impact Statement}: TURRET OS establishes a new paradigm for counter-espionage defense by demonstrating that multi-layered provenance knowledge graphs can simultaneously resolve out-of-distribution adversarial evasion and deliver court-admissible digital forensic evidence at zero operational delay.
